% Generated draft: requires mathematical and source review.
\documentclass[11pt]{article}
\usepackage[margin=25mm]{geometry}
\usepackage{fontspec}
\setmainfont{Latin Modern Roman}
\usepackage{amsmath,amssymb,mathtools,unicode-math}
\setmathfont{Latin Modern Math}
\usepackage{enumitem,needspace,xurl,hyperref}
\setlength{\parindent}{0pt}
\setlength{\parskip}{5pt}
\setlength{\emergencystretch}{3em}
\newcommand{\N}{\mathbb{N}}
\newcommand{\Z}{\mathbb{Z}}
\newcommand{\Q}{\mathbb{Q}}
\newcommand{\R}{\mathbb{R}}
\newcommand{\C}{\mathbb{C}}
\newcommand{\F}{\mathbb{F}}
\newcommand{\A}{\mathbb{A}}
\newcommand{\PP}{\mathbb{P}}
\newcommand{\E}{\mathbb{E}}
\newcommand{\eps}{\varepsilon}
\newcommand{\dd}{\,\mathrm{d}}
\begin{document}
\section*{UnsolvedMath ID: 30002707}
\textbf{Source problem code:} catalogue entry
\subsection{Definitions and mathematical statement}
A semialgebraic set is a finite Boolean combination of sets defined by polynomial equalities and inequalities; a function is semialgebraic when its graph is semialgebraic. The coefficient field for these polynomial descriptions is not specified in the supplied record. This omission matters when interpreting the question as a statement about Kontsevich--Zagier periods.

The supplied question is as follows. Let \(X_1,X_2\subset\R^n\) be semialgebraic sets, and let \(f_1,f_2\) be semialgebraic functions on their respective domains for which the integrals converge. If
\[
\int_{X_1} f_1=\int_{X_2} f_2,
\]
must there exist a finite sequence of integral identities transforming the first displayed integral into the second, with each identity justified by additivity, a semialgebraic change of variables, or Stokes' formula?

Here a finite sequence means that each successive integral expression is obtained from the preceding expression by an allowed rule, possibly involving sums of integrals. This explains the intended transformation question; it does not supply a missing formal calculus of rules. The record does not specify the coefficient field of admissible intermediate data, whether convergence must be absolute, the precise meaning of improper integration, or the regularity, orientation, and boundary hypotheses required for the stated rules. Those qualifications remain unresolved. In particular, no restriction to rational or algebraic coefficients is silently imposed, and the supplied formulation cannot without further clarification be identified with a fully specified arithmetic period conjecture.
\subsection{Short English statement}
If two convergent integrals of semialgebraic functions over semialgebraic sets have equal values, must additivity, semialgebraic substitutions, and Stokes' formula connect their expressions in finitely many steps? The source presents this as the real Kontsevich--Zagier period question, but leaves the coefficient field and exact admissibility conditions unspecified.
\subsection{Sources}
\begin{itemize}
\item Supplied local registry record: \texttt{data/unsolvedmath/problems.json}, permanent UnsolvedMath ID \(30002707\), exact source problem code \texttt{catalogue entry}; title: ``Kontsevich--Zagier Equivalence of Real Period Integrals.'' The record gives proposed year 2014 and names no proposer.
\item Supplied report citation: \emph{Valuation Theory and Its Applications (2015)}, Oberwolfach Reports, DOI 10.4171/owr/2014/49, pp.\ 2757--2823. \url{https://doi.org/10.4171/owr/2014/49}. The supplied background describes Oberwolfach Reports as documenting mathematical workshops and their open-problem discussions.
\item Supplied literature source: ``Kontsevich--Zagier period conjecture overview.'' \url{https://numdam.org/articles/10.5802/jtnb.1204/}.
\item Supplied literature source: ``Oberwolfach Report 49/2014.'' \url{https://doi.org/10.4171/owr/2014/49}.
\item Literature assessment dated exactly \texttt{2026-08-21}: the supplied curation labels the problem open and states that the implication from equal real integral values to equivalence under the standard rules remains open; it also states that the report's \(p\)-adic discussion does not settle the real claim. This is supplied dated curation, not an independently verified literature assessment. The background requests independent verification of time-sensitive status and cited claims.
\item The registry labels the statement \texttt{corrected\_verified} and reports that it was checked against the official Oberwolfach report, preserving its hypotheses, notation, and open target. These are supplied verification claims; no source consultation or independent mathematical verification is asserted here.
\end{itemize}
\end{document}
